\section{Προσεγγιστικά Μοντέλα Αεροδυναμικής (Surrogate Models)}
\label{sec:bg_surrogates}

Τα προσεγγιστικά μοντέλα (Surrogate Models / Metamodels)~\cite{jin2011surrogate} αντικαθιστούν τις υπολογιστικά ακριβές προσομοιώσεις με ταχείς εκτιμητές μηχανικής μάθησης. Στο πλαίσιο της βελτιστοποίησης αεροδυναμικής, οι πλέον διαδεδομένες προσεγγίσεις περιλαμβάνουν:
\begin{enumerate}
    \item \textbf{Διαδικασίες Gauss (Gaussian Process Regression - Kriging):}~\cite{rasmussen2006gaussian} Παρέχουν μη-παραμετρική παλινδρόμηση και παράλληλη εκτίμηση της αβεβαιότητας $\sigma^2(\bm{x})$ της πρόβλεψης:
    \begin{equation}
    \hat{y}(\bm{x}) = \bm{k}^\top (\bm{K} + \sigma_n^2 \bm{I})^{-1} \bm{y}
    \end{equation}
    \item \textbf{Πολυεπίπεδα Νευρωνικά Δίκτυα (MLP):} Ταχύτατη εκτέλεση συμπερασμού (inference) με δυνατότητα συνεχούς προσαρμογής μέσω στοχαστικής καθόδου κλίσης (SGD/Adam).
    \item \textbf{Random Forests (RF):} Σύνολα δέντρων αποφάσεων ανθεκτικά σε ακραίες τιμές και διακυμάνσεις κλίμακας.
    \item \textbf{k-Πλησιέστεροι Γείτονες (k-NN):} Μη-παραμετρική τοπική παρεμβολή βασισμένη σε Ευκλείδειες αποστάσεις.
\end{enumerate}
